\documentclass[11pt]{article}
\usepackage{amsmath,amsthm,amssymb}
\usepackage[margin=1in]{geometry}
\usepackage{booktabs}
\usepackage{array}
\usepackage{stmaryrd}
\usepackage[colorlinks=true,linkcolor=blue,citecolor=blue,urlcolor=blue]{hyperref}

\newtheorem{theorem}{Theorem}
\newtheorem{proposition}[theorem]{Proposition}
\newtheorem{corollary}[theorem]{Corollary}
\newtheorem{lemma}[theorem]{Lemma}
\theoremstyle{definition}
\newtheorem{definition}[theorem]{Definition}
\newtheorem{remark}[theorem]{Remark}
\newtheorem{example}[theorem]{Example}

\newcommand{\SPoly}{\mathsf{SPoly}}
\newcommand{\PolyN}{\mathsf{Poly}_{\mathbb{N}}}
\newcommand{\Poly}{\mathsf{Poly}}
\newcommand{\Cont}{\mathsf{Cont}}
\newcommand{\Set}{\mathsf{Set}}
\newcommand{\CRing}{\mathsf{CRing}}
\newcommand{\N}{\mathbb{N}}
\newcommand{\Z}{\mathbb{Z}}
\newcommand{\DD}{\mathbb{D}}
\newcommand{\sem}[1]{\llbracket #1 \rrbracket}
\newcommand{\id}{\mathrm{id}}

\title{The container derivative has no fiberwise negation\\
\large A boundary for the Cartan calculus of polynomial functors}
\author{MacBeth\thanks{Kodamai. Source, this note, and the Lean certification live at
\texttt{github.com/scotmacbeth/work-in-progress}; see the content commit recorded on this
page. For Rick, as the Cartan-boundary sequel to the container-derivative uniqueness note.}\\
\small Kodamai}
\date{5 October 2026\\[2pt]
\small content commit \texttt{5a33e34}}

\begin{document}
\maketitle

\begin{abstract}
The one-hole derivative of Abbott, Altenkirch, Ghani and McBride is now known to be the
differential combinator of a Cartesian differential category of polynomial functors, and
to be the unique nontrivial representable first-order tangent structure on the
finite-support ones; existence and uniqueness are settled. The natural next step is the
differential geometry built on top: the Aintablian--Blohmann program generates the full
Cartan calculus --- exterior derivative, contraction, Lie derivative, and the
Lie--Rinehart assembly of forms --- from a tangent category equipped with a \emph{scalar
ring object}. We show this step fails for containers, and locate exactly why. The additive
bundle of the dual-number tangent structure $T=(-)\lhd D$ is a bundle of commutative
monoids whose fibers are free $\N$-modules on holes; they have no additive inverses. Hence
$(\SPoly,\partial)$ is not a Rosick\'y tangent category and carries no scalar ring object,
and the full Cartan calculus does not descend. The obstruction is not an absence of ring
objects --- $\Z$ and $\Z/n$ do occur in $(\Poly,\times)$ as constant functors --- but the
impossibility of a \emph{compatible scalar action} on $T$, which would negate tangent
vectors. What survives is the subtraction-free fragment that uses only $+$, $0$, and
$\N$-scaling. The rig $\N$ is the knife, exactly as it was for uniqueness.
\end{abstract}

\section{Introduction}\label{sec:intro}

\paragraph{From a derivative to a differential geometry.}
Containers \cite{AAGM2003,GambinoKock2013} have a derivative. The one-hole-context
operation $\partial$ of Abbott, Altenkirch, Ghani and McBride sends a container to the
container of its one-hole contexts, and it obeys the laws of a derivative: a sum rule, a
product rule, and a chain rule. Two companion notes make this precise and sharp. The first
\cite{MacBethTangent} shows that the category $\SPoly$ of finitary polynomial functors is a
Cartesian differential category whose differential combinator \emph{is} $\partial$, with
the attached tangent bundle computed by the oldest trick in differential algebra ---
substitute a dual number and read off the $\varepsilon$-coefficient,
\[
   T(p)\ =\ p(y+\varepsilon v)\bmod \varepsilon^2\ =\ p(y)+\varepsilon\, v\cdot\partial p .
\]
The second \cite{MacBethUniqueness} shows this tangent structure is essentially forced:
among representable first-order tangent structures $T=(-)\lhd D$ on \emph{finite-support}
polynomial functors there are exactly two, the trivial one and $\partial$. Existence and
uniqueness are banked.

So containers have a canonical first-order differential calculus: directional derivatives,
a tangent bundle, the chain rule. The obvious next ambition --- the ambition behind this
whole program --- is the \emph{second}-order and tensorial apparatus of differential
geometry: differential forms, the exterior derivative $d$, contraction $\iota_X$ with a
vector field, the Lie derivative $L_X$, and the Chevalley--Eilenberg / Lie--Rinehart
assembly that packages a tangent structure into a de Rham-like complex. If containers have
a derivative, do they have a Cartan calculus?

\paragraph{The input the Cartan calculus demands.}
There is now a clean categorical answer to what that apparatus needs. Aintablian and
Blohmann \cite{AintablianBlohmann2026} construct the entire Cartan calculus on an arbitrary
Cartesian tangent category, provided it carries one extra piece of data: a \emph{scalar
ring object} $R$ --- a commutative ring object acting on each tangent bundle so that the
bundle becomes an $R$-module internally. From $(R,\kappa)$ they generate $d$, $\iota$, $L$,
and the Lie--Rinehart algebra of vector fields; their sequel \cite{AintablianBlohmann2026b}
builds a commutative differential graded algebra of cubical forms on the same input. The
scalar ring object is the hinge. Everything hangs on whether $\SPoly$ has one.

The decisive structural fact about a scalar ring object is elementary and classical: an
$R$-module has additive inverses, because multiplying a vector by $-1\in R$ negates it. So
a scalar ring object forces the tangent bundle to be a bundle of \emph{abelian groups} ---
it forces the tangent category to be \emph{Rosick\'y}, or \emph{additive with negatives}.
The whole question collapses to one line:

\begin{center}
\emph{Do tangent vectors on containers have negatives?}
\end{center}

\paragraph{The answer, and the result.}
They do not. A tangent vector on a container is, concretely, an $\N$-valued multiplicity
attached to each hole --- how many infinitesimal steps to take in each hole-direction ---
and these multiplicities add in $\N$. There is no subtraction, because there is nothing to
subtract: positions are elements of a set, and a set has no additive inverses. The tangent
fiber over a shape with $k$ holes is the free $\N$-module $\N^k$, a commutative monoid in
which only $0$ is invertible.

\begin{theorem}[Boundary theorem]\label{thm:main-intro}
The additive bundle of the tangent structure $T=(-)\lhd D$, $D=y+\varepsilon v$, on
finite-support polynomial functors is a bundle of commutative monoids that does not admit
fiberwise additive inverses. Consequently $(\SPoly,\partial)$ is not a Rosick\'y tangent
category and carries no scalar ring object in the sense of
\cite{AintablianBlohmann2026}; the Aintablian--Blohmann Cartan calculus does not descend to
$(\SPoly,\partial)$. At most the subtraction-free, $\N$-scalar fragment survives.
\end{theorem}

\paragraph{Method.}
The proof has two faces, which settle the same obstruction twice. The airtight one
(\S\ref{sec:transport}) is a \emph{transport} argument. The counting functor
$N:\SPoly\to\PolyN$, which replaces each container by the ordinary $\N$-coefficient
polynomial counting its shapes and positions, is a morphism of Cartesian tangent
categories \cite{MacBethTangent}. A morphism of tangent categories carries the additive
bundle --- and hence any candidate negation --- to the target. But in $\PolyN$ a negation
of the dual-number tangent structure would be an $\N$-coefficient polynomial $g(x,v)$ with
$v+g(x,v)=0$, forcing $g=-v$, which has a negative coefficient. No such polynomial exists.
The argument uses only that $N$ \emph{preserves} the finitely many operations in the
negation axiom, so no exotic negation can escape downward. The concrete face
(\S\ref{sec:fibers}) computes the tangent fibers directly in the bundle presentation that
the Cartan program actually consumes, and exhibits the free-$\N$-module obstruction on the
minimal witnesses $p=y$ (the line) and $p=y^2$.

\paragraph{The honest knife.}
It would be a mistake to read this as ``$\Poly$ has no ring objects.'' It does: every
ordinary commutative ring $k$ --- in particular $\Z$ and $\Z/n$ --- sits inside
$(\Poly,\times)$ as a constant functor, and the constants form a full, product-preserving
copy of $\CRing$. The obstruction is sharper and more interesting: no ring object can
\emph{act} on the dual-number tangent bundle $T=(-)\lhd D$ by a scalar multiplication
compatible with the tangent structure, because any such action would, through
multiplication by $-1$, negate tangent vectors --- and the fibers admit no negation. The
knife is the compatible action, not the ring. This is the same subtraction-free knife that
forced the uniqueness converse to be reproved rig-internally
\cite{MacBethUniqueness}: \emph{rig-versus-ring is the dividing line}, now as a theorem on
the differential-geometric side.

\paragraph{Outline.}
Section~\ref{sec:prelim} fixes the categories, the tangent structure, and the notion of
scalar ring object. Section~\ref{sec:lemmaA} reproves, from the module axioms, that a
scalar ring object forces negatives. Section~\ref{sec:transport} proves the main
obstruction by transport through the counting functor;
Section~\ref{sec:fibers} exhibits the same obstruction as an explicit fiber computation in
the Cartan-facing presentation. Section~\ref{sec:consequences} draws the consequences and
delimits precisely what fails and what survives. Section~\ref{sec:lean} records the scope
of the machine-checked certification. Section~\ref{sec:conclusion} points to the
group-completion sequel.

\paragraph{Scope.}
Everything here is first order, as are $\partial$ itself, the existence and uniqueness
notes, and the fibrational tangent theory of \cite{CCGZ2024} on which the framing rests.

\section{Preliminaries}\label{sec:prelim}

We recall only what the proof uses.

\paragraph{Containers and polynomial functors.}
A \emph{container} $S\lhd a$ is a set of \emph{shapes} $S$ together with, for each $s\in
S$, a \emph{position set} $a_s$; it is \emph{finite-support} (equivalently
\emph{finitary}) when every $a_s$ is finite. Its extension is the polynomial functor
\[
   \sem{S\lhd a}(X)\ =\ \sum_{s\in S} X^{a_s},
\]
and only the cardinalities $|a_s|$ matter up to isomorphism. We write $\SPoly$ for the
category of finitary polynomial functors; $\lhd$ denotes substitution (composition of
extensions), with unit the identity functor $y$. The AAGM derivative $\partial$ sends
$S\lhd a$ to the container of one-hole contexts, $\partial\sem{S\lhd a}(X)=\sum_{s}\sum_{h
\in a_s} X^{a_s\setminus\{h\}}$; on a single shape with $k$ positions,
$\partial(y^k)=k\,y^{k-1}$, the holes being the $k$ ways to delete a position
\cite{AAGM2003}.

\paragraph{The dual-number tangent structure.}
Let $D=y+\varepsilon v$ be the dual-number infinitesimal object, with
$\varepsilon^2=0$ and infinitesimal part $M=y$. The tangent bundle is substitution into
$D$,
\[
   T(p)\ =\ p\lhd D\ =\ p(y+\varepsilon v)\bmod\varepsilon^2\ =\ p + \varepsilon\, v\cdot
   \partial p ,
\]
with projection $p_{\mathrm{proj}}:Tp\to p$ (set $\varepsilon=0$) and zero section
$0:p\to Tp$ (set $v=0$) \cite{MacBethTangent}. We use the standard tangent-category
vocabulary of \cite{CockettCruttwell2014}: the \emph{additive bundle} is the
commutative-monoid structure on each fiber of $p_{\mathrm{proj}}$, given by a fiberwise
addition $+:T_2p\to Tp$ (on the pullback $T_2p=Tp\times_p Tp$ of two tangent vectors over a
common base point) and the zero section. A tangent category is \emph{Rosick\'y}, or
\emph{additive with negatives}, when each such fiber is an abelian group: there is a bundle
endomorphism $\sigma:T\Rightarrow T$ over the projection with
\begin{equation}
   +\ \circ\ \langle \sigma,\ \id\rangle\ =\ 0\circ p_{\mathrm{proj}}
   \qquad\text{as maps } T(X)\to T(X). \tag{$\dagger$}\label{eq:neg}
\end{equation}
We call such a $\sigma$ a \emph{negation}.

\paragraph{The scalar ring object.}
We take verbatim the input the Cartan calculus consumes.

\begin{definition}[scalar ring object; {\cite[\S4]{AintablianBlohmann2026}}]
\label{def:scalar-ring}
In a Cartesian tangent category $(\mathcal C, T)$, a \emph{scalar ring object} is a
commutative unital ring object $(R,\hat+,\hat\cdot,\hat 0,\hat 1,\hat-)$ for the Cartesian
product $\times$, together with a natural \emph{scalar multiplication} $\kappa_X:R\times
T(X)\to T(X)$ making each additive tangent bundle $p_{\mathrm{proj}}:T(X)\to X$ an
$R$-module in the slice $\mathcal C/X$: unital and associative over $\hat\cdot$, bilinear
with respect to $\hat+$ and the tangent addition $+_X$, fiber-preserving
($p_{\mathrm{proj}}\circ\kappa_X=p_{\mathrm{proj}}\circ\pi_2$), and compatible with the
vertical lift and the canonical flip. This is the data from which
\cite{AintablianBlohmann2026} generate $d$, $\iota$, $L$ and the Lie--Rinehart algebra of
vector fields.
\end{definition}

We rely on \cite{AintablianBlohmann2026} only for this \emph{definition of the input}; the
consequence we use (Lemma~\ref{lem:A}) is reproved below.

\paragraph{The counting functor.}
The counting functor $N:\SPoly\to\PolyN$ sends a container to the $\N$-coefficient
polynomial that counts its data: $N(\sum_s y^{a_s}) = \sum_s x^{|a_s|}\in\N[x]$, and
multivariably by counting positions of each colour. Here $\PolyN$ is the Cartesian
differential (hence tangent) category of polynomials over the rig $\N$ --- objects $\N$,
morphisms tuples of $\N$-coefficient polynomials, differential the Jacobian directional
derivative, tangent structure $T(k)=2k$, $T(h)=\langle h\circ\pi_1, D[h]\rangle$. By
\cite[Lemmas 3--5]{MacBethTangent}, $N$ preserves composition, finite products and
projections, the left-additive $+$ and $0$, the product of functors, and the differential
($N\partial=\partial N$); it is therefore a morphism of Cartesian tangent categories in the
sense of \cite{CockettCruttwell2014}. We will need only the preservation, not the
faithfulness.

\section{A scalar ring object forces negatives}\label{sec:lemmaA}

The reduction is a single internalised fact from module theory.

\begin{lemma}\label{lem:A}
If $(R,\kappa)$ is a scalar ring object on a Cartesian tangent category $(\mathcal C,T)$,
then
\[
   \sigma_X\ :=\ \kappa_X\circ\big\langle\, (\hat-\circ\hat1\circ\, !_R),\ \id_{T X}\,\big\rangle
   \ :\ T(X)\longrightarrow T(X)
\]
--- scalar multiplication by $-1:=\hat-\hat1\in R$ --- is a negation. Hence $(\mathcal
C,T)$ is Rosick\'y.
\end{lemma}

\begin{proof}
This is the standard fact that in any module over a ring, multiplication by $-1$ is the
additive inverse, read internally. By the $R$-module axioms on the bundle
$p_{\mathrm{proj}}:TX\to X$ --- bilinearity of $\kappa$ over $\hat+$, unitality, and the
ring unit/zero laws --- for any tangent vector $t$ over a base point $x$,
\[
   t\ +_X\ (-1)\cdot t\ =\ 1\cdot t\ +_X\ (-1)\cdot t\ =\ \big(\hat1\,\hat+\,(\hat-\hat1)\big)\cdot t
   \ =\ \hat0\cdot t\ =\ 0_X(x),
\]
using unitality $1\cdot t=t$, bilinearity $(r\hat+s)\cdot t=r\cdot t\ +_X\ s\cdot t$, the
ring law $\hat1\,\hat+\,(\hat-\hat1)=\hat0$, and $\hat0\cdot t=0_X$. The last equality is
itself forced by bilinearity: $\hat0\cdot t=(\hat0\,\hat+\,\hat0)\cdot t=\hat0\cdot t\ +_X\
\hat0\cdot t$, and cancelling (fibers are commutative monoids) gives $\hat0\cdot t=0_X$.
Naturality of $\kappa$ and the fiber-preservation law $p_{\mathrm{proj}}\circ\kappa_X=
p_{\mathrm{proj}}\circ\pi_2$ give $p_{\mathrm{proj}}\circ\sigma_X=p_{\mathrm{proj}}$, so
$\sigma$ is a bundle endomorphism satisfying \eqref{eq:neg}.
\end{proof}

Lemma~\ref{lem:A} is \cite[Prop.~4.4]{AintablianBlohmann2026} in substance, reproved here
from the module axioms so that the boundary theorem depends only on the
\emph{definition} of a scalar ring object, not on trusting any derived statement. Its
contrapositive is the engine of the paper:

\begin{quote}
\emph{If some additive-bundle fiber of $(\SPoly,T)$ is a commutative monoid that is not a
group, then $(\SPoly,T)$ has no scalar ring object, and the Cartan calculus is obstructed
at its very input.}
\end{quote}

The whole program thus reduces to one decidable question: \emph{is the additive bundle of
$T=(-)\lhd D$ a bundle of groups?} The next two sections answer no.

\section{The additive bundle has no fiberwise negation: transport}\label{sec:transport}

We work in the Cartesian-differential presentation of $\SPoly$, where the counting functor
$N$ is a genuine functor and carries the full tangent structure to $\PolyN$. The slogan:
\emph{a negation upstairs would push down to a negation in $\PolyN$, and $\PolyN$ has
none.}

\begin{proposition}\label{prop:transport}
The container tangent structure $(\SPoly,T)$ admits no negation $\sigma$.
\end{proposition}

\begin{proof}
Suppose $\sigma:T\Rightarrow T$ were a negation, i.e.\ a bundle endomorphism satisfying
\eqref{eq:neg}. Apply $N$. Because $N$ preserves composition, finite products and
projections, the identity, the zero section, the projection, the additive pullback $T_2$,
and the fiberwise addition $+$ --- all the operations appearing in \eqref{eq:neg}, and the
pairing $\langle-,-\rangle$ into $T_2$ is a universal construction from products and the
projection, which $N$ also preserves --- equation \eqref{eq:neg} is carried to
\[
   +\ \circ\ \langle N(\sigma),\ \id\rangle\ =\ 0\circ p_{\mathrm{proj}}
   \qquad\text{in } \PolyN .
\]
Thus $N(\sigma)$ is a negation for the dual-number tangent structure of $\PolyN$.

But $\PolyN$ has none. Take the object $1$; then $T(1)=1\times 1$ has two coordinates
$(x,v)$, base and tangent, with projection $p_{\mathrm{proj}}=\pi_x$, zero section $x\mapsto
(x,0)$, and fiberwise addition $(x,v_1,v_2)\mapsto(x,v_1+v_2)$. A bundle endomorphism over
$p_{\mathrm{proj}}$ is a map $(x,v)\mapsto(x,g(x,v))$ for some $\N$-coefficient polynomial
$g$ --- morphisms of $\PolyN$ are exactly $\N$-polynomials --- and the negation law
\eqref{eq:neg} reads
\[
   v + g(x,v)\ =\ 0 \qquad\text{identically in } \N[x,v].
\]
This forces $g=-v$, which is not an $\N$-polynomial: it has a negative coefficient, and
already at $v=1$ it demands $g(x,1)=-1$, a value no $\N$-polynomial attains (they are
nonnegative on $\N$). Contradiction. Hence no negation exists, and the additive bundle
$p_{\mathrm{proj}}:T\Rightarrow\id$ is a bundle of commutative monoids that are not groups.
\end{proof}

\begin{remark}[why this cannot be dodged]
The argument uses only that $N$ \emph{preserves} the finitely many operations occurring in
\eqref{eq:neg} --- not that it is faithful. Any candidate negation whatsoever, however
exotic, is transported downward to a negation of $\PolyN$ and killed there. Structurally:
$\PolyN$ is a left-additive category whose hom-monoids $(\N[x],+)$ are not groups, so it is
not an additive-with-negatives tangent category; $N$ pushes any negation down, so
$(\SPoly,T)$ cannot carry one either.
\end{remark}

\section{The same obstruction in the fibers}\label{sec:fibers}

The Cartan program consumes the tangent bundle in its \emph{bundle} presentation: objects
are the polynomial functors $p$ themselves --- the ``spaces'' --- and $Tp=p\lhd D$ is the
tangent bundle of the space $p$, with additive bundle $p_{\mathrm{proj}}:Tp\to p$. The
obstruction of \S\ref{sec:transport} appears here as an explicit fiber computation on
small, named spaces. The two presentations share one additive bundle (the dual-number
$\varepsilon v$-part, with the same coproduct $\N$-module addition), so the verdict is the
same; the point of this section is to make the free-$\N$-module structure tangible.

For $p=\sum_s y^{a_s}$ the bundle $Tp=p\lhd D = p+\varepsilon v\cdot\partial p$ has base
part $p$ and tangent part
\[
   \varepsilon v\cdot\partial p\ =\ \varepsilon v\cdot\sum_s\sum_{h\in a_s} y^{a_s\setminus\{h\}},
\]
one tangent direction per hole $h\in a_s$ over each shape $s$. The additive bundle of $Tp$
is $p\lhd(\text{additive bundle of }D)$, where $D$ is the commutative-monoid object with
infinitesimal part $M=y$: the fiberwise addition $+:T_2p\to Tp$ is $p\lhd\nabla$ with
$\nabla:M\oplus M\to M$ the codiagonal, which collapses two infinitesimal copies
$\varepsilon v_1,\varepsilon v_2$ to $\varepsilon v$. Iterating $\nabla$ --- the
$\N$-module structure that every additive bundle carries by repeated fiber addition ---
makes the tangent data over a shape $s$ the \emph{free $\N$-module on the holes $a_s$}: a
tangent vector carries an $\N$-multiplicity per hole-direction, added hole by hole in $\N$.
There is no addition of position-values to appeal to; positions are elements of a set. The
counting polynomial realises this fiber as $(\N^{|a_s|},+,0)$, in agreement with
\S\ref{sec:transport}. A negation would make each fiber the free $\Z$-module instead.

\begin{example}[the line $p=y$; minimal witness]
One shape, one hole; $Ty=y+\varepsilon v=D$, fiber $\N$. The unit tangent $1\in\N$ has no
additive inverse: $1+c=0$ is unsolvable in $\N$. Already the tangent bundle of the
\emph{identity functor} has a non-group fiber.
\end{example}

\begin{example}[the first nonlinear space $p=y^2$]
$\partial(y^2)=2y$, so $T(y^2)=y^2+\varepsilon v\cdot 2y$; one shape, two holes
$\{h_1,h_2\}$, fiber $\N^2$. The element $(1,0)$ --- one unit of tangent along $h_1$ --- has
no inverse in $\N^2$.
\end{example}

\begin{example}[a mixed space $p=y+y^2$]
$\partial p=1+2y$. There are two shapes: $s_0$ with no holes (fiber $\N^0=0$, trivially a
group) and $s_1$ with two holes (fiber $\N^2$). The $s_1$-fiber already fails, so the
bundle is not a bundle of groups.
\end{example}

Each fiber is a genuine commutative monoid --- commutative, cancellative, with identity
$0$ --- but only $0$ is invertible. This is exactly the free $\N$-module structure; a group
would be the free $\Z$-module, and $\Z$ is one subtraction beyond $\N$.

\section{Consequences: the honest boundary}\label{sec:consequences}

\begin{corollary}[no scalar ring object]\label{cor:no-ring}
$(\SPoly,\ T=(-)\lhd D)$ carries no scalar ring object in the sense of
Definition~\ref{def:scalar-ring}.
\end{corollary}

\begin{proof}
Such an $R$ would, by Lemma~\ref{lem:A}, make $(\SPoly,T)$ Rosick\'y --- its additive
bundle a bundle of abelian groups. But that additive bundle has a fiber that is not a
group: the free $\N$-module $\N^{|a_s|}$ over any shape with a hole
(\S\ref{sec:fibers}), equivalently the $(\N,+)$ obstruction transported by $N$
(Proposition~\ref{prop:transport}). Two independent witnesses, one contradiction.
\end{proof}

\begin{remark}[the precise obstruction --- not ``no ring objects'']
Ring objects \emph{do} exist in $(\Poly,\times)$. Every ordinary commutative ring $k$ ---
in particular $\Z$ and $\Z/n$ --- embeds as a constant polynomial functor, and the
constant functors form a full, finite-product-preserving copy of $\CRing$ inside $\Poly$.
The obstruction in Corollary~\ref{cor:no-ring} is therefore \emph{not} the absence of a
ring object. It is that no ring object can \emph{act} on the dual-number tangent bundle
$T=(-)\lhd D$ by a scalar multiplication $\kappa$ compatible with the tangent structure:
any such action would, via $\kappa(\hat-\hat1,-)$, negate tangent vectors
(Lemma~\ref{lem:A}), and the fibers admit no negation
(Proposition~\ref{prop:transport}). The knife is the compatible action, not the ring.
\end{remark}

\paragraph{What does not descend.}
The full Aintablian--Blohmann package built on $(R,\kappa)$ --- the exterior derivative
$d$, the contraction $\iota_X$, the Lie derivative $L_X$, and the
Lie--Rinehart / Chevalley--Eilenberg assembly of forms, and a fortiori the cubical-forms
differential graded algebra of \cite{AintablianBlohmann2026b}, which additionally requires
$R$ to have no $2$-torsion --- does not descend to $(\SPoly,\partial)$. In particular the
Chevalley--Eilenberg differential, which needs signs, is the first casualty of the missing
$-1$.

\paragraph{What survives: the $\N$-scalar fragment.}
The additive bundle is nonetheless a genuine $\N$-module bundle: the fibers are free
$\N$-modules on holes, with a canonical $\N$-scaling $c\cdot(-)$ for every $c\in\N$,
realised by iterating the fiberwise addition. Every operation of the calculus that uses
only $+$, $0$, and $\N$-scaling --- not $-1$ --- survives on $(\SPoly,\partial)$. Pinning
down exactly which of $d$ and $\iota$ survive the loss of subtraction is the natural
follow-up; we flag it rather than claim it.

\paragraph{The grant frame.}
Existence and uniqueness are banked; this theorem \emph{delimits} the program honestly. The
differential geometry of containers is a genuine first-order tangent calculus ---
directional derivatives, a tangent bundle, the chain rule --- but not a full Cartan
calculus with forms and Lie derivatives. The rig $\N$ blocks the scalar ring. ``Rig-versus
-ring is the knife,'' already visible in the uniqueness converse
\cite{MacBethUniqueness}, becomes a theorem on the geometric side: the exact slide after
uniqueness.

\section{Scope of the certification}\label{sec:lean}

The arithmetic core of the obstruction is machine-checked in Lean~4 (Mathlib-free; axiom
profile \texttt{[propext, Quot.sound]} only, via \texttt{omega}). The formalised statements
are: that the negation law $v+g(x,v)=0$ has no solution over the rig $\N$ --- indeed no
function $\N\times\N\to\N$ solves it, a fortiori no $\N$-polynomial
(Proposition~\ref{prop:transport}); that the free $\N$-module $\N^k$ ($k\ge1$) admits no
fiberwise negation, is cancellative, and has only $0$ invertible; and the explicit fiber
witnesses for $p=y$ and $p=y^2$ (\S\ref{sec:fibers}). An auxiliary Python script checks the
dual-number identities $Tp=p+\varepsilon v\cdot\partial p$ and
$N(Tp)=N(p)(x+\varepsilon v)$ on a spread of examples.

We state the scope precisely and without inflation. The Lean development certifies the
\emph{arithmetic} obstruction --- the non-existence of a negation over $\N$ and the
free-$\N$-module structure of the fibers. It does \emph{not} formalise Lemma~\ref{lem:A}
(the module-algebra reduction, which is pen-and-paper) nor the functoriality of the
counting functor $N$ (established in the existence note, also pen-and-paper). This matches
the convention of the companion notes, whose finite oracles certify the arithmetic heart of
each result and leave the categorical scaffolding to prose.

\section{Conclusion and sequel}\label{sec:conclusion}

Containers have a derivative and a first-order tangent calculus, and no more: the dual
-number tangent bundle has no fiberwise negation, so there is no scalar ring object and the
Cartan calculus of forms and Lie derivatives does not descend. The single cause is the rig
$\N$ --- the subtraction-free arithmetic of positions-as-set-elements --- and the honest
statement is not that containers lack rings but that no ring can act on their tangent
vectors by a negation-inducing scalar.

This locates the one move that might lift the ceiling. The fibers are free $\N$-modules;
their obstruction is precisely the obstruction of $\N$ to being a ring, namely the missing
group completion. Fiberwise group completion $\N\to\Z$ turns each free $\N$-module into a
free $\Z$-module and each fiber into a group. The question the boundary theorem exposes ---
and the most interesting one it leaves open --- is whether this completion is
\emph{categorical}: does group-completing the tangent fibers land in a larger,
polynomial-like category of \emph{virtual} or $\Z$-containers that \emph{is} Rosick\'y, and
on which the full Cartan calculus therefore \emph{does} descend? That is the honest route to
a differential geometry of containers with forms, and it is the subject of the sequel.

The result is unpublished; a precedent to acknowledge is Cockett's 2012 talk
\cite{Cockett2012}, which first proposed that polynomial functors form a Cartesian
differential category. We know it only secondhand, through a community relay of the
unpublished slides, with Cruttwell's caveat that combinatorial species are more naturally a
tangent \emph{bicategory} than the $1$-categorical $\SPoly$ used here; the two framings must
not be conflated.

\begin{thebibliography}{9}

\bibitem{AAGM2003}
M.~Abbott, T.~Altenkirch, N.~Ghani, C.~McBride,
\emph{Derivatives of containers},
Typed Lambda Calculi and Applications (TLCA 2003), LNCS 2701, Springer, 2003, 16--30.

\bibitem{AhmanUustalu2016}
D.~Ahman, T.~Uustalu,
\emph{Directed containers as categories},
Proc.\ MSFP 2016, EPTCS 207, 2016, 89--98.

\bibitem{AintablianBlohmann2026}
A.~Aintablian, C.~Blohmann,
\emph{Cartan calculus in tangent categories},
preprint, arXiv:2607.11169 (2026).

\bibitem{AintablianBlohmann2026b}
A.~Aintablian, C.~Blohmann,
\emph{Cartan calculus of cubical forms in tangent categories},
preprint, arXiv:2609.05963 (2026).

\bibitem{BCS2009}
R.~F.~Blute, J.~R.~B.~Cockett, R.~A.~G.~Seely,
\emph{Cartesian differential categories},
Theory and Applications of Categories \textbf{22} (2009), no.~23, 622--672.

\bibitem{CCGZ2024}
M.~Capucci, G.~Cruttwell, N.~Ghani, F.~Zanasi,
\emph{A fibrational theory of first order differential structures},
preprint, arXiv:2409.05763 (2024).

\bibitem{Cockett2012}
J.~R.~B.~Cockett,
\emph{Can you Differentiate a Polynomial?},
slides, Foundational Methods in Computer Science (FMCS 2012), unpublished; known to the
author secondhand via a community relay.

\bibitem{CockettCruttwell2014}
J.~R.~B.~Cockett, G.~S.~H.~Cruttwell,
\emph{Differential structure, tangent structure, and SDG},
Applied Categorical Structures \textbf{22} (2014), no.~2, 331--417.

\bibitem{GambinoKock2013}
N.~Gambino, J.~Kock,
\emph{Polynomial functors and polynomial monads},
Math.\ Proc.\ Cambridge Philos.\ Soc.\ \textbf{154} (2013), no.~1, 153--192.

\bibitem{LanfranchiLemay2025}
M.~Lanfranchi, J.-S.~P.~Lemay,
\emph{Representable tangent structures for affine schemes},
preprint, arXiv:2505.09080 (2025).

\bibitem{MacBethTangent}
MacBeth,
\emph{The container derivative is a tangent structure: a fibrational positioning}
(Kodamai internal note, 2026).

\bibitem{MacBethUniqueness}
MacBeth,
\emph{A uniqueness theorem for the container derivative: $\partial$ is the only nontrivial
representable tangent structure on finite-support polynomial functors}
(Kodamai internal note, 2026).

\end{thebibliography}

\end{document}
